\input{../tex-inputs/latex-korrekturansicht-vorspann}

\section[Arthur Schnitzler an Theodor Herzl, 21. 6. 1893]{L03902 Arthur Schnitzler an Theodor Herzl, 21. 6. 1893}
\nopagebreak\mylabel{L03902v}
\rehead{ }\normalsize\beginnumbering\briefempfaengerindex{Herzl, Theodor@\textsc{Herzl, Theodor}!zzzSchnitzler, Arthur@\emph{von Arthur Schnitzler}!1893-06-211@{21. 6. 1893}|(be}
\toendnotes[C]{\smallbreak\pagebreak[2]}
\correspDesc{Versand  durch Arthur Schnitzler am 21. 6. 1893 in Wien
\newline{}Erhalt  durch Theodor Herzl im Zeitraum [22. 6. 1893
                  – 29. 6. 1893] in Paris}\toendnotes[C]{\smallbreak}
\Standort{Jerusalem, Central Zionist Archives, H1:1924-7.}
\physDesc{Brief, 2 Blätter, 8 Seiten, 2649 Zeichen (Briefpapier mit Trauerrand)
\newline{}Handschrift: schwarze Tinte, deutsche Kurrent
\newline{}Ordnung: mit Bleistift von unbekannter Hand innerhalb das Konvoluts paginiert:
                                    »25«–»28« }
\buchAbdrucke{\weitereDrucke{Arthur Schnitzler: \emph{Briefe 1875–1912}. Herausgegeben von Therese Nickl und Heinrich Schnitzler. Frankfurt am Main: \emph{S. Fischer} 1981, S. 207–209.} }\toendnotes[C]{\smallbreak}
\pstart{}{\pb}Verehrter Freund,\pend\vspace{0.5em}
\pstart
           hoffentlich iſt alles bei Ihnen wohl und die Maſern haben Ihre \textcolor{blue}{älteſte}\pwindex{Hüft, Pauline 29.\,3.\,1890 – 8.\,9.\,1930@\textsc{Hüft, Pauline} (29.\,3.\,1890 – 8.\,9.\,1930)|pwv}{}\ledrightnote{{$\rightarrow$}\emph{\textcolor{blue}{Pauline Hüft}}} verſchont. Es wird also wohl nach
                  \textcolor{pink}{Wien}\oindex{Wien@\textbf{Wien}, \emph{Verwaltungsgebiet}|pw}{}\ledrightnote{\textcolor{pink}{Wien}} gefahren? – Wie geſagt, ich wiederhole
               meine Bitte, daſs Sie mich irgendwie verſtändigen, wie lange Sie dableiben \textsc{etc etc.} – Es iſt nicht unmöglich, daß ich Anfang Juli auf 10–14 Tage mit meiner \label{K_L03902-1v}\edtext{armen \textcolor{blue}{Mama}\pwindex{Schnitzler, Louise 8.\,7.\,1840 Kőszeg – 9.\,9.\,1911 Wien@\textsc{Schnitzler, Louise} (8.\,7.\,1840 Kőszeg – 9.\,9.\,1911 Wien)|pwv}{}\ledrightnote{{$\rightarrow$}\emph{\textcolor{blue}{Louise Schnitzler}}}}{\lemma{\textnormal{\emph{armen Mama}}}\Cendnote{\textnormal{Sie war seit dem 2. 5. 1893 Witwe.}}}\label{K_L03902-1}{ }\label{K_L03902-2v}\edtext{nach \textcolor{pink}{Iſchl}\oindex{Bad Ischl@\textbf{Bad Ischl}|pw}{}\ledrightnote{\textcolor{pink}{Bad Ischl}}}{\lemma{\textnormal{\emph{nach Iſchl}}}\Cendnote{\textnormal{\textcolor{blue}{Schnitzler} reiste am 2. 7. 1893 alleine nach \textcolor{pink}{Ischl}\oindex{Bad Ischl@\textbf{Bad Ischl}|pwk}. Am 16. 7. 1893 kam er
                  wieder in \textcolor{pink}{Wien}\oindex{Wien@\textbf{Wien}, \emph{Verwaltungsgebiet}|pwk} an.}}}\label{K_L03902-2} gehe; im übrigen wäre ein biſſel Gebirgs{\pb}luft für mich nicht gerade überflüſſig, da ich ſeit den
               qualvollen Aufregungen des Frühjahrs an einem großentheils nervöſen Huſten leide, den
               ich nicht anbringe. Im übrigen habe ich mich jetzt auf den Sport geworfen u. fahre
               ſeit ein \label{K_L03902-3v}\edtext{paar Tagen auf dem \textsc{Bicycle}}{\lemma{\textnormal{\emph{paar … Bicycle}}}\Cendnote{\textnormal{Vgl. A. S.: \emph{Tagebuch}, 13. 6. 1893.}}}\label{K_L03902-3}. Ich ſchreibe dieſe
               Zeilen mit ſteifem Arm und ſteifem Bein – iſt auch das letztere zu bemerken?–\pend
           
\pstart
           {\pb}– Ihren \textcolor{green}{Prinzen aus
                  Genieland}\pwindex{Herzl, Theodor 2.\,5.\,1860 Budapest – 3.\,7.\,1904 Edlach@\textsc{Herzl, Theodor} (2.\,5.\,1860 Budapest – 3.\,7.\,1904 Edlach), \emph{Schriftsteller, Journalist}!Prinzen aus Genieland. Lustspiel in 4 Akten@\strich\emph{Prinzen aus Genieland. Lustspiel in 4 Akten}|pw}{}\ledrightnote{\textcolor{green}{Prinzen aus Genieland. Lustspiel in 4 Akten}} hab ich hier im \textcolor{pink}{Carltheater}\oindex{Wien@\textbf{Wien}!II., Leopoldstadt@\textbf{II., Leopoldstadt}!Carl-Theater@\textbf{Carl-Theater}, \emph{Theater}|pw}{}\ledrightnote{\textcolor{pink}{Carl-Theater}}{ }\textcolor{violet}{geſehen}\eventindex{Carl-Theater@\textbf{Carl-Theater}!Uraufführung von Prinzen in Genieland, 20.11.1891@Uraufführung von Prinzen in Genieland, 20.11.1891|pwv}{}\ledrightnote{{$\rightarrow$}\emph{\textcolor{violet}{Uraufführung von Prinzen in Genieland, 20.11.1891}}} – we{\geminationn} ſie nicht ganz um\substVorne{}\textsuperscript{lege}\substDazwischen{}gebr\substHinten{}acht worden ſind, ſo ſind nicht die Mörder, ſondern die Prinzen ſelbſt \introOben{}dran\introOben{} ſchuld geweſen; denn es iſt im ganzen miſerabel geſpielt
               worden. Mir war das \textcolor{green}{Stück}\pwindex{Herzl, Theodor 2.\,5.\,1860 Budapest – 3.\,7.\,1904 Edlach@\textsc{Herzl, Theodor} (2.\,5.\,1860 Budapest – 3.\,7.\,1904 Edlach), \emph{Schriftsteller, Journalist}!Prinzen aus Genieland. Lustspiel in 4 Akten@\strich\emph{Prinzen aus Genieland. Lustspiel in 4 Akten}|pwv}{}\ledrightnote{{$\rightarrow$}\emph{\textcolor{green}{Prinzen aus Genieland. Lustspiel in 4 Akten}}} ſehr
               ſympathiſch; es lag wie ein Duft von 1840 drüber; es gehört vielleicht zu denjenigen
               Ihrer dramatiſchen Sachen, in {\pb}denen am meiſten Frühling
               iſt. Allerdings sind beträchtliche Ungleichheiten drin, und was mir am ärgerlichſten
               daran war, – ſo weit mich heut meine Eri{\geminationn}erg nicht
               täuſcht – war die zu pathetiſche u abſichtliche Manier, in welchem plötzlich die
               Grundidee (im \textcolor{green}{3. Akt}\pwindex{Herzl, Theodor 2.\,5.\,1860 Budapest – 3.\,7.\,1904 Edlach@\textsc{Herzl, Theodor} (2.\,5.\,1860 Budapest – 3.\,7.\,1904 Edlach), \emph{Schriftsteller, Journalist}!Prinzen aus Genieland. Lustspiel in 4 Akten@\strich\emph{Prinzen aus Genieland. Lustspiel in 4 Akten}|pwv}{}\ledrightnote{{$\rightarrow$}\emph{\textcolor{green}{Prinzen aus Genieland. Lustspiel in 4 Akten}}} glaub
               ich) \substVorne{}\textsuperscript{\textcolor{gray}{d}}\substDazwischen{}a\substHinten{}usgeſprochen wird, ſtatt daſs das ganze \textcolor{green}{Stück}\pwindex{Herzl, Theodor 2.\,5.\,1860 Budapest – 3.\,7.\,1904 Edlach@\textsc{Herzl, Theodor} (2.\,5.\,1860 Budapest – 3.\,7.\,1904 Edlach), \emph{Schriftsteller, Journalist}!Prinzen aus Genieland. Lustspiel in 4 Akten@\strich\emph{Prinzen aus Genieland. Lustspiel in 4 Akten}|pwv}{}\ledrightnote{{$\rightarrow$}\emph{\textcolor{green}{Prinzen aus Genieland. Lustspiel in 4 Akten}}} im Fortſchreiten ſelbſt und in ſeinen {\pb}Charakteren jene Grundidee ausſpricht. – Sehr deutlich iſt
               mir eine treffliche Charge \textsc{\textcolor{blue}{Knaack}\pwindex{Knaack, Wilhelm 13.\,2.\,1829 Rostock – 29.\,10.\,1894@\textsc{Knaack, Wilhelm} (13.\,2.\,1829 Rostock – 29.\,10.\,1894), \emph{Schauspieler}|pw}{}\ledrightnote{\textcolor{blue}{Wilhelm Knaack}}’}s im Gedächtnis. – Da{\geminationn} die hübſche kleinbürgerliche Scene im \textcolor{green}{2. Akt}\pwindex{Herzl, Theodor 2.\,5.\,1860 Budapest – 3.\,7.\,1904 Edlach@\textsc{Herzl, Theodor} (2.\,5.\,1860 Budapest – 3.\,7.\,1904 Edlach), \emph{Schriftsteller, Journalist}!Prinzen aus Genieland. Lustspiel in 4 Akten@\strich\emph{Prinzen aus Genieland. Lustspiel in 4 Akten}|pwv}{}\ledrightnote{{$\rightarrow$}\emph{\textcolor{green}{Prinzen aus Genieland. Lustspiel in 4 Akten}}}, in der eine Violine vorko{\geminationm}t. Da{\geminationn} die Liebesſcene,
               die von Herrn \textsc{\textcolor{blue}{Franker}\pwindex{Franker, Heinrich *~1857@\textsc{Franker, Heinrich} (*~1857)|pw}{}\ledrightnote{\textcolor{blue}{Heinrich Franker}}} u Frl. \textsc{\textcolor{blue}{Ernst}\pwindex{Ernst, Carla 15.\,5.\,1867 Wien – 15.\,6.\,1925 ebd.@\textsc{Ernst, Carla} (15.\,5.\,1867 Wien – 15.\,6.\,1925 ebd.), \emph{Schauspielerin}|pw}{}\ledrightnote{\textcolor{blue}{Carla Ernst}}} gar nicht übel gegeben wurde. Unangenehm war der Oberprinz – Herr \textsc{\textcolor{blue}{Lenor}\pwindex{Lenor, Robert v. 8.\,1.\,1855 Wien – 29.\,6.\,1900 Innsbruck@\textsc{Lenor, Robert v.} (8.\,1.\,1855 Wien – 29.\,6.\,1900 Innsbruck), \emph{Schauspieler, Verwaltungsbeamter}|pw}{}\ledrightnote{\textcolor{blue}{Robert v. Lenor}}}, der seine Rolle {\pb}mit einer ſchnodderigen
               Liebenswürdigkeit gab, welche mir im Ohr und im Herzen wehthat. – Ich hatte damals
               den lebhaften Wunſch, das \textcolor{green}{Stück}\pwindex{Herzl, Theodor 2.\,5.\,1860 Budapest – 3.\,7.\,1904 Edlach@\textsc{Herzl, Theodor} (2.\,5.\,1860 Budapest – 3.\,7.\,1904 Edlach), \emph{Schriftsteller, Journalist}!Prinzen aus Genieland. Lustspiel in 4 Akten@\strich\emph{Prinzen aus Genieland. Lustspiel in 4 Akten}|pwv}{}\ledrightnote{{$\rightarrow$}\emph{\textcolor{green}{Prinzen aus Genieland. Lustspiel in 4 Akten}}} wo anders zu ſehen, u mit mir wünſchten ſich einige Vernünftige, daſs
               das Luſtſpiel nicht am \textcolor{brown}{Volkstheater}\orgindex{Volkstheater@Volkstheater|pw}{}\ledrightnote{\textcolor{brown}{Volkstheater}} gegeben würde
               – \textsc{resp.} dß Sie’s nicht dem \textcolor{brown}{V.th.}\orgindex{Volkstheater@Volkstheater|pw}{}\ledrightnote{\textcolor{brown}{Volkstheater}}{ }{\pb}eingereicht hätten. – Ihre Geschichten von \textcolor{blue}{Barnay}\pwindex{Barnay, Ludwig 11.\,2.\,1842 Budapest – 1.\,2.\,1924@\textsc{Barnay, Ludwig} (11.\,2.\,1842 Budapest – 1.\,2.\,1924), \emph{Schriftsteller, Schauspieler, Theaterdirektor}|pw}{}\ledrightnote{\textcolor{blue}{Ludwig Barnay}}, vom \textcolor{green}{Flüchtling}\pwindex{Herzl, Theodor 2.\,5.\,1860 Budapest – 3.\,7.\,1904 Edlach@\textsc{Herzl, Theodor} (2.\,5.\,1860 Budapest – 3.\,7.\,1904 Edlach), \emph{Schriftsteller, Journalist}!Flüchtling. Lustspiel in einem Aufzug@\strich\emph{Der Flüchtling. Lustspiel in einem Aufzug}|pw}{}\ledrightnote{\textcolor{green}{Der Flüchtling. Lustspiel in einem Aufzug}} u. ſo w. haben mich intereſſirt u. gerührt. Jawohl gerührt – denn
               ich finde ſolche Abenteuer des Ehrgeizes und des Talents um die Ehrlichkeit rührender
               als hungrige Mägen und manches verliebte Idyll. – Vom \textcolor{brown}{Volkstheater}\orgindex{Volkstheater@Volkstheater|pw}{}\ledrightnote{\textcolor{brown}{Volkstheater}} müſſen Sie mir erzählen.\pend
           
\pstart
           {\pb}– Der Brief da trifft \substVorne{}\textsuperscript{ſ}\substDazwischen{}S\substHinten{}ie, lieber Freund, wohl noch in \textcolor{pink}{Paris}\oindex{Paris@\textbf{Paris}, \emph{Hauptstadt}|pw}{}\ledrightnote{\textcolor{pink}{Paris}}.
               Nehmen Sie noch mein Glückwunsch zum \textcolor{blue}{Gretherl}\pwindex{Neumann, Margarethe 20.\,5.\,1893 Paris – 15.\,3.\,1943 Konzentrationslager Theresienstadt@\textsc{Neumann, Margarethe} (20.\,5.\,1893 Paris – 15.\,3.\,1943 Konzentrationslager Theresienstadt)|pw}{}\ledrightnote{\textcolor{blue}{Margarethe Neumann}}
               entgegen und beſtellen Sie dieſelben auch Ihrer Frau \textcolor{blue}{Gemahlin}\pwindex{Herzl, Julie 1.\,2.\,1868 Budapest – 10.\,11.\,1907 Wien@\textsc{Herzl, Julie} (1.\,2.\,1868 Budapest – 10.\,11.\,1907 Wien)|pwv}{}\ledrightnote{{$\rightarrow$}\emph{\textcolor{blue}{Julie Herzl}}}, die ſich vielleicht noch meiner erinnert.–\pend
           
\pstart
           Alſo auf baldiges Wiederſehen,{\\[\baselineskip]}Ihr herzlich ergebener{\\[\baselineskip]}\spacefill\mbox{ArthSchnitzler}\pend
           \leftskip=0em{}
\pstart
           21. 6. 93.\pend
           \selectlanguage{ngerman}\endnumbering\briefempfaengerindex{Herzl, Theodor@\textsc{Herzl, Theodor}!zzzSchnitzler, Arthur@\emph{von Arthur Schnitzler}!1893-06-211@{21. 6. 1893}|)be}\mylabel{L03902h}  \newcommand{\dateiname}{L03902}\newcommand{\titel}{Arthur Schnitzler an Theodor Herzl, 21. 6. 1893}\newcommand{\editorInnen}{Herausgegeben von Jahnke, SelmaMüller, Martin Anton}\input{../tex-inputs/latex-korrekturansicht-abspann}
